% !TeX root = main.tex
\newpage
\section{Future Modification and Cleanup}

\subsection{Boundary-Space Lighting}

\subsubsection{Schedule}

The first boundary-space lighting milestone is complete for
\texttt{LightingBall}.  Python deposits persistent color on existing boundary
proxies, and Vulkan \texttt{ParticleLighting} uses the \texttt{FPML.*} shader
family plus \texttt{BoundaryLightRecord} records to produce the same visible
result.

\subsubsection{Validated Form}

The first implementation should prove:
\[
    \textrm{photon hit}
    \rightarrow
    \textrm{persistent boundary proxy color}
    \rightarrow
    \textrm{kernel-spread surface color}.
\]

This uses existing boundary particles as the first persistent proxies.  Do not
restore the retired \texttt{BoundaryLightSample}/\texttt{BoundaryLightState}
arrays, visible surface-cell grid, eye-driven brightness, or per-frame
clear/filter path.

\subsubsection{Deferred Work}

The next Vulkan rendering refinement is to use \texttt{gl\_PointCoord} in the
lighting fragment shader so each boundary point renders as a soft patch
instead of a hard point.  Directional specular, additive energy accumulation,
temporal filtering, decay, and a separate compact Vulkan proxy table remain
deferred.

\subsection{Vulkan Test-File View Defaults}

Move the Vulkan initial view defaults that are scene/data specific out of the
application configuration and into the generated \texttt{.tst} file.  In
particular, the current setting:
\begin{verbatim}
initial_view = { rRotZ = 0.0;rRotY=0.0;rRotX=0.0;zoom=0.1;};
\end{verbatim}
should be generated with the test data and read from the \texttt{.tst} file,
not from the Vulkan app configuration file.  This keeps \texttt{view_center},
point size, and initial zoom controlled by the same scene-specific data source.

\subsection{Photon-Direction Specular Lighting}

\subsection{Boundary-Light Washout Cap}

Add an optional render control for LightingBall boundary lighting:
\begin{verbatim}
boundary_light_washout_limit = 0.75;
\end{verbatim}

This would let a lit yellow sector become pale yellow without becoming pure
white.  The center of a dense photon patch could still wash out strongly, but
the final mix toward white would be capped by the configured limit.

\subsubsection{Schedule}

Defer photon-direction specular lighting until the scalar boundary-light path
is proven in both Python and Vulkan:
\[
    \textrm{photon hit}
    \rightarrow
    \textrm{boundary sample energy}
    \rightarrow
    \textrm{filtered lit sector}.
\]

\subsubsection{Target Idea}

Later, boundary light samples may accumulate the incoming photon direction as
well as deposited energy.  A coherent set of photon directions could then
produce a specular highlight from the simulation itself rather than from a
configured analytic light vector.  This should be a second-stage extension,
not part of the first LightingBall memory contract.
